\clearpage
\VerdierPhysicalPageAnchor{0188}%
\begingroup\small\setlength{\parskip}{0pt}\setlength{\emergencystretch}{1.5em}
\setlength{\abovedisplayskip}{4pt plus 2pt minus 2pt}%
\setlength{\belowdisplayskip}{4pt plus 2pt minus 2pt}%

\VerdierSemanticAnchor{ch3d-2-1-3-properties-cont}
% Anomalie diplomatique de l'autorité : « un morphismes » ; correction : « un morphisme ».
un morphisme élémentaire. Les morphismes élémentaires sont stables par
changement de base. Soit :

\VerdierSemanticAnchor{ch3d-2-1-3-bi-infinite-chain}
\[
\cdots\longrightarrow X_n\xrightarrow{u_n}X_{n-1}
\xrightarrow{u_{n-1}}X_{n-2}\longrightarrow\cdots
\]
une suite de morphismes de complexes, illimitée dans les deux sens, telle
que pour tout entier $n$, le morphisme $u_n$ soit un morphisme élémentaire
de hauteur $-n$. Cette suite admet dans
$\operatorname{comp}(\mathcal A)$ une limite projective et inductive (cette
suite induit sur les composants de degré donné une suite stationnaire dans
les deux sens).

\VerdierSemanticAnchor{ch3d-2-1-4}
\medskip\noindent\textbf{Lemme 2.1.4.}\quad
\textit{Soit $m\colon X\longrightarrow Y$ un quasi-isomorphisme de
$\operatorname{comp}(\mathcal A)$. Il existe une suite de morphismes :}
\[
u_n\colon X_n\longrightarrow X_{n-1},\qquad n\in\mathbb Z,
\]
\noindent\textit{telle que $u_n$ soit, pour tout entier $n$, un morphisme
élémentaire de hauteur $-n$, et des morphismes :}

\VerdierSemanticAnchor{ch3d-2-1-4-limits}
\[
\varinjlim_n X_n\xleftarrow{\sim}Y,
\qquad
\varprojlim_n X_n\xleftarrow{v}X,
\qquad v\text{ homotopisme (chap. I, 2.5.11)},
\]
\noindent\textit{tels que le diagramme :}

\VerdierSemanticAnchor{ch3d-2-1-4-square}
\begin{center}
\begin{tikzcd}[column sep=5.6em,row sep=3em]
X \arrow[r,"v"] \arrow[d,"m"']
  & \displaystyle\varprojlim_n X_n \arrow[d] \\
Y \arrow[r,"{\sim}"'] & \displaystyle\varinjlim_n X_n
\end{tikzcd}
\end{center}
\noindent\textit{soit commutatif à homotopie près.}

\VerdierSemanticAnchor{ch3d-2-1-4-proof}
Tout d'abord, remarquons que pour tout morphisme
$m\colon X\longrightarrow Y$ de $\operatorname{comp}(\mathcal A)$, il existe
un épimorphisme $m'\colon X'\longrightarrow Y$ et un diagramme commutatif à
homotopie près :

\VerdierSemanticAnchor{ch3d-2-1-4-epimorphism-diagram}
\begin{center}
\begin{tikzcd}[column sep=3.5em,row sep=3.2em]
X \arrow[rr,"v"] \arrow[dr,"m"'] && X' \arrow[dl,"{m'}"] \\
& Y &
\end{tikzcd}\qquad,
\end{center}
\endgroup
